% Additive research note for R7-01 / D0033, 2026-09-17.
% This file does not replace the existing manuscript.

\section{Arcsine tensorization and first-order sharpness inside stationary Markov chains}

We retain the notation of the main proof.  For a Markov transition kernel
$K(x,\mathrm dy)$ on a standard Borel space, define its robust Hilbertian
contraction coefficient
\[
 \kappa(K):=\sup_{\lambda}\rho(U,V),
 \qquad U\sim\lambda,\quad V\mid U=x\sim K(x,\cdot),
\]
where the supremum is over all input laws supported on the full-measure set of
rows used by the chain (degenerate inputs contribute zero).  For the chain let
\[
 \kappa(n):=\sup_{t-s\ge n}\kappa(K_{s,t}).
\]
The Dobrushin/chi-square lemma proved in the preceding note gives
\[
 \kappa(K)^2\le \delta(K),
 \qquad
 \kappa(n)^2\le \Delta(n)\le \min\{1,2\phi(n)\}.
\tag{1}
\]

\subsection{An arcsine improvement of the block-path bound}

Define
\[
 \Psi(r):=\sin\!\left(\min\{2\arcsin r,\pi/2\}\right),\qquad 0\le r\le1.
\]
Thus
\[
 \Psi(r)=2r\sqrt{1-r^2}\quad(0\le r\le 2^{-1/2}),
 \qquad \Psi(r)=1\quad(r\ge2^{-1/2}).
\tag{2}
\]

\paragraph{Theorem.}
For every lag $n$,
\[
 \boxed{\rho^*(n)\le \Psi(\kappa(n)).}
\tag{3}
\]
Consequently, whenever $\phi(n)\le1/4$,
\[
 \boxed{\rho^*(n)\le
  2\sqrt{2\phi(n)\bigl(1-2\phi(n)\bigr)}.}
\tag{4}
\]
In particular the source implication $\phi(n)\to0\Rightarrow\rho^*(n)\to0$
remains valid, while (4) strictly improves the previously recorded
$2\sqrt{2\phi(n)}$ bound whenever $0<\phi(n)<1/4$.

\paragraph{Proof.}
Fix finite disjoint index sets $S,T$ with $\operatorname{dist}(S,T)\ge n$.
Order $S\cup T$ chronologically and merge consecutive selected indices of the
same color into maximal monochromatic blocks.  Write the $S$-blocks in
chronological order as $U_i$ and the $T$-blocks as $V_j$.  After possibly
shifting one index, a block $U_i$ is adjacent in the chronological block path
only to $V_i$ and $V_{i-1}$.

Use Peyre's minimal-hypothesis formulation of the tensorization theorem: the
subjective correlation of $(U_i,V_j)$ is taken after conditioning only on
$U_{i'}$, $i'<i$, and $V_{j'}$, $j'<j$.  If the two blocks are nonadjacent,
this conditioning contains an intervening Markov separator; past and future
are therefore conditionally independent and the subjective correlation is
zero.  If the blocks are adjacent, let $s$ be the last time in the earlier
block and $t$ the first time in the later block.  The prescribed conditioning
is entirely earlier than this transition.  It can change the input law of
$X_s$, but it does not change $K_{s,t}$, and the conditioned input remains
supported on the same row-good set.  By definition of $\kappa$ and Markov data
processing,
\[
 \rho(U_i,V_j\mid\hbox{prescribed prefixes})
 \le \kappa(K_{s,t})\le\kappa(n).
\tag{5}
\]

Pad the two finite block sequences by deterministic variables and index them
by $\mathbb Z$.  Peyre's ``$\mathbb Z$ against $\mathbb Z$'' theorem
(Theorem 3.3.10, with the minimal hypotheses from Section 3.4(a)) applies with
\[
 \varepsilon(0)=\varepsilon(-1)=\kappa(n),\qquad
 \varepsilon(z)=0\quad(z\notin\{0,-1\}).
\]
It gives
\[
 \arcsin\rho(\sigma(U_i:i),\sigma(V_j:j))
 \le \min\{2\arcsin\kappa(n),\pi/2\},
\]
which is (3).  Taking the supremum over finite interlacings and then the usual
finite-coordinate $L^2$ approximation proves (3) for $\rho^*(n)$.
Equation (4) follows from (1), monotonicity of $\Psi$ on $[0,2^{-1/2}]$, and
$\Delta(n)\le2\phi(n)$.  \qed

Peyre's Section 3.5 records optimality of the arcsine theorem under its abstract
prefix-correlation hypotheses.  Hence any improvement beyond (3) must exploit
additional compatibility specific to Markov transition kernels, rather than
only the two-edge path of subjective correlation bounds.

\subsection{A stationary two-state family forcing the path factor two}

The factor two in the weak-dependence expansion of (3) is already forced by a
stationary two-state Markov chain.  Let $0<a<1$ and let $(X_k)_{k\in\mathbb Z}$
take values in $\{-1,+1\}$ with stationary law $(1/2,1/2)$ and transition rule
\[
 \mathbb P(X_{k+1}=X_k)=\frac{1+a}{2},\qquad
 \mathbb P(X_{k+1}=-X_k)=\frac{1-a}{2}.
\tag{6}
\]
Then
\[
 \mathbb E[X_{k+h}\mid X_k]=a^hX_k,
 \qquad \mathbb E[X_iX_j]=a^{|i-j|}.
\tag{7}
\]
Fix $n\ge1$ and put $r=a^n$.  The $n$-step channel is the binary symmetric
channel with parameter $r$.  Its robust maximal correlation is exactly $r$.
Indeed, for an input law with mean $m$, the only centered nonconstant input
direction is $X-m$, and direct calculation gives
\[
 \rho(U,V)^2=
 \frac{r^2(1-m^2)}{1-r^2m^2}\le r^2,
\tag{8}
\]
with equality at the stationary input $m=0$.  Also
\[
 \boxed{\phi(n)=\frac r2.}
\tag{9}
\]
For (9), conditioning the past can only alter the posterior law of $X_0$;
after $n$ steps the resulting bias is multiplied by $r$, so total variation
from stationarity is at most $r/2$.  Equality is attained by conditioning on
$X_0=+1$ and testing $X_n=+1$.

Now choose the interlacing
\[
 S_N=\{0,2n,\ldots,2Nn\},\qquad
 T_N=\{n,3n,\ldots,(2N-1)n\},
\]
and the admissible centered test functions
\[
 F_N=\sum_{i=0}^{N}X_{2in},\qquad
 G_N=\sum_{j=0}^{N-1}X_{(2j+1)n}.
\]
From (7), counting pairs by their separation yields
\[
 \frac{\operatorname{Var}(F_N)}{N}
 \longrightarrow 1+2\sum_{h\ge1}r^{2h}
 =\frac{1+r^2}{1-r^2},
\tag{10}
\]
and the same limit for $\operatorname{Var}(G_N)/N$.  Likewise each bulk odd
site sees even sites at odd distances, so
\[
 \frac{\operatorname{Cov}(F_N,G_N)}{N}
 \longrightarrow \sum_{k\in\mathbb Z}r^{|2k-1|}
 =\frac{2r}{1-r^2}.
\tag{11}
\]
Boundary terms in (10)--(11) are $O(1)$ because the corresponding geometric
series are summable.  Therefore
\[
 \lim_{N\to\infty}\operatorname{Corr}(F_N,G_N)
 =\boxed{\frac{2r}{1+r^2}},
\tag{12}
\]
and hence
\[
 \boxed{\rho^*(n)\ge\frac{2a^n}{1+a^{2n}}
       =\frac{4\phi(n)}{1+4\phi(n)^2}.}
\tag{13}
\]

On the other hand, applying (3) with the exact robust local coefficient
$\kappa(n)=r$ gives, for $r\le2^{-1/2}$,
\[
 \rho^*(n)\le2r\sqrt{1-r^2}.
\tag{14}
\]
Thus as $r\downarrow0$,
\[
 2r-2r^3+O(r^5)
 \le \rho^*(n)
 \le 2r-r^3+O(r^5),
\]
so
\[
 \boxed{\rho^*(n)=2r+O(r^3)}
 \qquad(r\downarrow0).
\tag{15}
\]
In particular the coefficient $2$ in the weak-edge expansion of the path
tensorization step cannot be replaced by any smaller universal coefficient,
even within stationary irreducible aperiodic two-state Markov chains.

This sharpness statement concerns the global path step as a function of the
robust local HGR coefficient $r$.  It does not show that the universal
Dobrushin conversion $\kappa^2\le\Delta\le2\phi$ is simultaneously sharp on
the same Markov family; indeed here $\kappa=r$ and $\Delta=r=2\phi$, so the
Dobrushin-to-HGR square-root bound is deliberately loose.  The remaining
constant question is therefore localized to compatibility between local
Dobrushin extremizers (erasure-type channels) and long stationary Markov
interlacings.
